Let us prove that (i) implies (ii), (iii), and (iv). Thus suppose that @@M0@@ is bounded by an invertible sheaf on an @@M1@@. Then @@M2@@ satisfies (iv) by EGA II 4.6.12. Suppose @@M3@@ flat over @@M4@@. Then the polynomial in @@M5@@ and @@M6@@, @@M7@@, depends only on the connected component of @@M8@@ (EGA III 7.9.4). Consequently, @@M9@@ satisfies (ii) and, taking 2.1 into account, (iii).

Next, note that the family of classes of the @@M10@@ (resp. @@M11@@) is bounded, hence satisfies (iii) (resp. (ii)).

Suppose that @@M12@@ satisfies (iv). Then @@M13@@, hence @@M14@@, and likewise @@M15@@. Moreover, @@M16@@, hence @@M17@@. Thus @@M18@@ satisfies (iii).

Finally, let us prove that (ii) (resp. (iii)) implies (i). By 2.10 and 2.3, one may replace @@M19@@ by a multiple\,; therefore suppose @@M20@@ very ample. Evidently one may decompose @@M21@@ into pieces\,; thus suppose that there exists a linear section @@M22@@ of @@M23@@ which is integral over @@M24@@, with fibers of dimension @@M25@@. By 3.8, one may replace @@M26@@ by @@M27@@, (taking account of 1.7 (resp. the projection formula), which implies that condition (ii) (resp. (iii)) on @@M28@@ is equivalent to (ii) (resp. (iii)) on @@M29@@)\,; therefore suppose @@M30@@. If @@M31@@ is @@M32@@, one may replace @@M33@@ by the family of classes of the @@M34@@ by 2.3 and 2.10\,; thus suppose @@M35@@.

If, in the Hilbert polynomials @@M36@@, @@M37@@ remains bounded and @@M38@@ remains bounded below, then the three coefficients remain bounded by 2.3 and 2.7\,; thus @@M39@@ is bounded by 2.11\,; hence (ii) implies
